% !TEX root = ../../../main.tex
% 来源：中学数学实验教材·第一册（代数）
% 章节：有理数系 / 有理数的意义
% 原始文件：1.tex，行 1196-1213
% 类型：tikz
% ---
	\begin{tikzpicture}[>=latex]
	\node (A) at (0,0) {有理数};
	\node (B1) at (2,1.5) {整数};
	\node (B2) at (2,-1.5) {分数};
	\node (B11) at (4,2.5)[right] {正整数（自然数）};
	\node (B12) at (4,1.5)[right] {零};
	\node (B13) at (4,.5)[right] {负整数};
	\node (B21) at (4,-.5)[right] {正分数};
	\node (B22) at (4,-2.5)[right] {负分数};
	\draw [->](A)--(B1);
	\draw [->](A)--(B2);
	\draw [->](B1)--(B11);
	\draw [->](B1)--(B12);
	\draw [->](B1)--(B13);
	\draw [->](B2)--(B21);
	\draw [->](B2)--(B22);
	
	\end{tikzpicture}
